\section{4. Sinus, cosinus og tangens}

\kompetence{Regning} \kompetence{Modellering} \kompetence{Hjaelpemidler}

\begin{figure}[ht]
\centering
\begin{tikzpicture}[>=Stealth,scale=2.5]
  % Koordinatsystem
  \draw[->] (-0.2,0) -- (1.3,0) node[below] {$x$};
  \draw[->] (0,-0.2) -- (0,1.3) node[left] {$y$};

  % Trekant
  \coordinate (A) at (0,0);
  \coordinate (B) at (1,0);
  \coordinate (C) at (1,0.75);

  % Trekant-sider
  \draw[very thick] (A) -- (B) node[midway,below] {$a$} -- (C) node[midway,right] {$b$} -- cycle;
  \draw[very thick,gxu-accent] (A) -- (C) node[midway,left] {$c$};

  % Vinkel
  \node at (0.2,0.1) {$\theta$};

  % Ret vinkel-markør
  \draw (0.92,0) -- (0.92,0.08) -- (1,0.08);

  % Hypotenuse-navn flyttet
  \node[anchor=south east] at (0.35,0.26) {$c$};
\end{tikzpicture}
\qquad
\begin{tikzpicture}[>=Stealth,scale=2.5]
  \pgfmathsetmacro{\cx}{cos(36.87)}
  \pgfmathsetmacro{\cy}{sin(36.87)}
  % Koordinatsystem
  \draw[->] (-0.2,0) -- (1.3,0) node[below] {$x$};
  \draw[->] (0,-0.2) -- (0,1.3) node[left] {$y$};
  \draw[thick] (0,0) circle (1);

  % Trekant: ret vinkel ved O, hypotenuse = radius til P
  \coordinate (P) at (\cx,\cy);
  \coordinate (B) at (\cx,0);
  \draw[thick] (0,0) -- (B) node[midway,below] {$1$} -- (P) -- cycle;
  \draw[very thick,gxu-accent] (0,0) -- (P) node[midway,left] {$\tan\theta$};

  % Ret vinkel-markør ved O
  \draw (0.08,0) -- (0.08,0.08) -- (0,0.08);

  % Radius = 1 label
  \node[anchor=north west] at (0.1,0.05) {$c=1$};

  % Vinkel-buе med pil
  \draw[->,thick,gxu-accent] (0.25,0) arc[start angle=0,end angle=36.87,radius=0.25];

  % Tekst
  \node[anchor=south west] at (0.55,0.10) {$a/b = \tan\theta$};
  \node[anchor=south west] at (0.55,0.26) {$b = \cos\theta$};
  \node[anchor=south west] at (0.55,0.42) {$a = \sin\theta$};

  % Navn P
  \node[anchor=south west] at (\cx+0.04,\cy) {$P=(\cos\theta,\sin\theta)$};
\end{tikzpicture}
\caption{Venstte: sinus, cosinus og tangens defineret fra en retvinklet trekant. Højre: de samme forhold aflæst på enhedscirklen.}
\label{fig:trigtrekant}
\end{figure}

For en retvinklet trekant med vinkel $\theta$ gælder:

\eq{4.1}{\sin\theta=\frac{a}{c}\qquad\cos\theta=\frac{b}{c}\qquad\tan\theta=\frac{a}{b}}

hvor $a$ er kateten over for $\theta$, $b$ er kateten hos $\theta$, og $c$ er hypotenusen.

\note{Vinklen $\theta$ er i grader.}

\begin{prompt}
I en retvinklet trekant er $a=5$ og $c=13$. Find $\sin\theta$, $\cos\theta$ og $\tan\theta$.
\end{prompt}
\begin{solutionc}
Først $b=\sqrt{13^{2}-5^{2}}=\sqrt{169-25}=\sqrt{144}=12$.\\
Derefter: $\sin\theta=5/13$, $\cos\theta=12/13$, $\tan\theta=5/12$.
\end{solutionc}
